\documentclass{article}
\usepackage[utf8]{inputenc}
\usepackage[T1]{fontenc}
\usepackage{geometry}
\geometry{
    a4paper,
    left=2.5cm,
    right=2.5cm,
    top=2.5cm,
    bottom=2.5cm
}
\usepackage{graphicx}
\usepackage{float}
\usepackage{hyperref}
\usepackage[english]{babel}
\usepackage{amsmath}
\usepackage{amssymb}
\usepackage{array}
\usepackage{booktabs}
\usepackage{caption}
\usepackage{subcaption}
\usepackage{enumitem}
\usepackage{fancyhdr}
\usepackage{titlesec}
\usepackage{url}
\usepackage{textcomp}
\usepackage{xcolor}
\usepackage{listings}
\usepackage{multirow}
\usepackage{longtable}
\usepackage{tabularx}
\usepackage{times}

% Set up page style
\pagestyle{fancy}
\fancyhf{}
\rfoot{\thepage}
\renewcommand{\headrulewidth}{0pt}

% Title formatting
\titleformat{\chapter}[display]{\normalfont\bfseries}{\chaptertitle\thechapter}{0pt}{\Huge}
\titlespacing*{\chapter}{0pt}{50pt}{40pt}

% Document info
\title{Template Document}
\author{Academic Support Agent}
\date{\today}

\begin{document}
\maketitle

\begin{center}
\textbf{XI International Conference on Computational Methods in Marine Engineering}

\bigskip
{\Large \textbf{Title of the Abstract}} \\
\bigskip

First A. Author$^1$, Second B. Author$^2$ and Third C. Author$^{1,*}$
  
$^1$ International Center for Numerical Methods in Engineering, \\
Universidad Polit cnica de Catalu a, Campus Norte UPC, 08034 Barcelona, Spain.

$^2$ School of Engineering, University of Edinburgh \\
The King s Buildings, Robert Stevenson Road, Sanderson Building Edinburgh, EH9 3FB, UK.

$^*$ Corresponding.Author@cimne.upc.edu.

\bigskip
{\Large \textbf{ABSTRACT}} \\
\bigskip
\end{center}

Either American or British spelling can be used, but one of them must be used consistently throughout the manuscript. Follow the Shorter Oxford English Dictionary. Manuscripts that have been carefully proofread will be considered for review, while manuscripts with unclear English phrases will be returned to the Corresponding Author without being reviewed. Pages must be numbered. The maximum number of pages is eight. Papers should feature a single-paragraph abstract of no more than 250 words. 

%%%%%%%%%%%%%%%%%%%%%%%%%%%%%%%%%%%%%%%%%%%%%%%%%%%%%%%%%%%%%%%%%%%%%%%%%%%%%%%%%%%%%%%%%%%%%%%

\textbf{Keywords:} list 6 to 10 keywords that could be used for an internet search by other researchers, e.g. lifting surface; hydrofoil hydrodynamics; wingsail; sail aerodynamics; leading-edge separation; impulse theory; lifting-line theory.

%%%%%%%%%%%%%%%%%%%%%%%%%%%%%%%%%%%%%%%%%%%%%%%%%%%%%%%%%%%%%%%%%%%%%%%%%%%%%%%%%%%%%%%%%%%%%%%

\section{INTRODUCTION} \label{sec:intro}
The paper should be structured in different sections. The recommended, but not compulsory list of sections is: Introduction, Method, Results, Conclusions, Acknowledgements (if any), and References. 
The first numbered section is the Introduction, and the last numbered section is the Conclusions. Headings are numbered, all capitalized, and left aligned. Sub-headings are numbered (1.1, 1.2, etc.), mixed case (initial capitals followed by lowercase on all words except articles, conjunctions, and prepositions, which should appear entirely in lowercase), and left aligned. Additional sub-headings are numbered (1.1.1, 1.1.2, etc.), mixed case (initial capitals followed by lower case on all words except articles, conjunctions, and prepositions, which should appear entirely in lower case), and left aligned.

\section{SYMBOLS} \label{sec:symbols}

Authors are kindly recommended to pay particular attention to the rigorous use of symbols to facilitate readability by a multidisciplinary readership. Each symbol should be in Italics only if it can take a numerical value and should be upright otherwise. For example, if the apparent wind angle appears in any equation, table, or plot, then a symbol in Italics such as $\beta$ must be defined at the first occurrence, and must be used in every equation, table and plot.  

Do not define symbols if they are not used in the rest of the manuscript. Define symbols in the Abstract only if used in the Abstract. When abbreviations are defined in the Abstract, they should also be defined again at the first occurrence in the main document, excluding the Abstract (e.g. in the Introduction). 

\subsection{Italics}

Two consecutive Italic symbols, either Roman or Greek letters, always identify the product of those. Therefore, do not use symbols with multiple letters, such as $TWA$, for true wind angle. Instead, when necessary, use subscripts, e.g. $\beta_\text{t}$. Dimensionless groups are a notable exception to this rule, as they can be defined by two letters, e.g. the Reynolds number is $Re$. Script font such as $\mathcal{G}$ and $\mathcal{R}$ can be used as an alternative to Italics when the same letter denotes a different quantity.

\subsection{Upright}

The upright font must always be used for operators, e.g. sin $\alpha$, log $y^+$, d$x$, etc.; constants such as i ($\sqrt{-1}$) and e (2.7183); physical units such as m, s, etc.; abbreviations such as w.r.t. (with respect to), CFD; and Latin words such as etc., et al., e.g., ad hoc. 
Bold Italic font (or bold sloping Greek) should be used for vectors, tensors and matrices, e.g. $\bm{u}=(u,v,w).$ However, consider using a different bold font for tensors and matrices such as $\bm{\mathcal{A,B,C,D}}$.

\subsection{Subscripts}

Some notations allow the subscripts to take a numerical value, for example, the indices of a tensor such as the shear stress tensor $T_{i,j}$, for which the first diagonal component is $T_{1,1}$. Therefore, the subscript should be in Italics or upright depending on whether these can take a numerical value, e.g. $T_{i,j}$ and $\beta_\text{t}$. Multiple subscripts should be separated by a comma if they are not multiplied, e.g. $T_{i,j}$. However, as subscripts are rarely multiplied between each other, it is acceptable to remove the comma to simplify the notation and write, for example, $T_{ij}$ instead of $T_{i,j}$. Note that the force along the coordinate $x$ can be indicated by $F_x$, where $x$ is in Italics because it refers to the symbol $x$. When the subscript indicates a word, it is always upright, e.g. $\beta_\mathrm{true}$. However, consider simplifying the notation by using the minimum number of letters or entirely avoid using the subscript if not necessary. 

\subsection{Products}

The product symbol ($\times$) should only be used to denote a cross product, or between numbers, e.g. $\bm{u}\times \bm{v}=\bm{w}, Re=5\times 10^5.$ The central point symbol ($\cdot$) should be used exclusively to denote a scalar product, e.g. $\bm{u}\cdot \bm{i}=0$. No symbol should be used to denote the product between two scalars, e.g. write $ab$ and not $a\cdot b$ nor $a\times b$.

\subsection{Uniqueness}

Each symbol must have a unique definition, for example, $\omega$ can be used either to denote vorticity or angular velocity, but not both. It is important to observe that the same symbol, upright or Italics, bold or lowercase, has different meanings. For example, $d$ might denote the diameter, while d is the differential operator. If $\bm{u}$ is defined as the velocity vector, then $u$ is intended as the magnitude of the vector, i.e. the speed. Therefore, it is critical that the symbols in the figures, in the equations and in the text are presented consistently in the appropriate format.  

\subsection{Punctuation}

Use punctuation as appropriate together with the equations, for example, `the lift coefficient of a sail with area $A$ in a uniform flow with freestream speed $u_\infty$ and density $\rho$ is
\begin{equation}\label{eq:CL}
    C_L\equiv \frac{L}{\frac{1}{2}\rho u_\infty^2 A},
\end{equation}\noindent
where $L$ is the lift.'  

\section{UNITS}\label{sec:units}

The International System of Units (SI), commonly known as the metric system, should be used throughout the document. Latex users may consider using the usepackage {\tt siunitx}. Units are always separated from the numbers by a non-breaking space, which is also used to separate different units. Multiple upright symbols non-broken by a space indicate a single unit, e.g. Pa (pascals), and $\mu$m (micrometre), while a non-breaking space must be used between units, e.g. m~s$^{-1}$. Note that units should be repeated for each entry in lists such as, for instance, \SIlist{1;2;3}{m}; and in ranges such as in the range \SIrange{1}{3}{m}. 

%%%%%%%%%%%%%%%%%%%%%%%%%%%%%%%%%%%%%%%%%%%%%%%%%%%%%%%%%%%%%%%%%%%%%%%%%%%%%%%%%%%%%%%%%%%%%%%

\section{ABBREVIATIONS}\label{sec:abbreviations}

All abbreviations must be defined at first use, even those deemed to be well known to the readership, such as 2D, 3D, CFD, etc. These are defined by writing the non-abbreviated terms using upper case for the letters used for the abbreviations when possible, followed by the abbreviation between brackets. For example, Two-Dimensional (2D) Computational Fluid Dynamics (CFD) simulations. 
Do not define abbreviations if they are not used in the rest of the manuscript. Define abbreviations in the Abstract only if used in the Abstract. When abbreviations are defined in the Abstract, they should also be defined again at the first occurrence in the main document, excluding the Abstract and the Nomenclature (e.g. in the Introduction). All abbreviations should be included in the Nomenclature, as in the following example. 

%%%%%%%%%%%%%%%%%%%%%%%%%%%%%%%%%%%%%%%%%%%%%%%%%%%%%%%%%%%%%%%%%%%%%%%%%%%%%%%%%%%%%%%%%%%%%%%

\section{FIGURES AND TABLES} \label{sec:figures}

Figures and tables should be embedded within the manuscript and appear as soon as possible after they have been cited in the text. Every figure and table must be cited in the text, must be numbered in the order that they appear, and must appear in the order in which they are first mentioned in the text. 

The words `figure,' `table,' `section,' and `equation' are non-capitalised only when they do not refer to a specific item in the paper. The words `figure' and `equation' are not abbreviated at the start of a sentence; otherwise, use Fig.~1, Figs.~1,2, Figs.~1,3-5,7, Eq.~2., Tab.~3, etc. Always use a non-breaking space before the number. Each figure and table should be accompanied by a single caption. Figure captions appear under the figure, while table cations appear above the table. 

Subfigures should be named as (a), (b), etc. Brackets around a non-capital letter should be used in the graphics, and in the figure caption, e.g. `... for (a) $Re=10^6$, and (b) $Re=10^8$'. Instead, in the main body of the manuscript, refer to the subfigure as fig.~\ref{fig:reattachmentlength}a, i.e. without brackets. 

When possible, include error bars in figures unless the uncertainty is clearly discussed in the text, and error bars would diminish the clarity of the figure. If the figure includes plots, the axes must be labelled with the corresponding symbols. For example, the horizontal axis presenting the time should be labelled as $t$ or $t$ [s] to highlight the units, not `time.' 

All authors must own or acquire the correct permissions and licences to reproduce figures. When possible, Authors are recommended to make their own figures. Authors do not need to acquire permission to reproduce data that is included in a new figure made by the Authors. 

The printable version of the manuscript in PDF format should be a maximum of 10 MB in size. Mindful of this limit, figures should be high resolution, e.g., 1200 dpi for lines, 300 dpi for halftone and colour in CMYK format, and 600 dpi for a mixture of lines and halftone. 

\begin{figure} 
	\begin{center}
    \includegraphics[width=0.6\textwidth]{./fig/ReattachmentLength.eps} 
 \caption{Nondimensional reattachment length versus the Reynolds number for values of the thickness-to-chord ratio spanning from 8 to 22.}
		\label{fig:reattachmentlength}
	\end{center}
\end{figure}

\begin{table}
  \caption{\textmd{Maximum lift and drag coefficients at different Reynolds numbers, and comparison with existing literature.}}
  \begin{center}
  \begin{tabular}{lccc} 
        & $Re=10^6$ & $Re=10^8$ & \cite{Viola2021OE} \\[3pt]
      $C_L$   & 1.1111 & 2.2222 & 3.3333\\
      $C_D$   & 4.4444 & 5.5555 & 6.6666
  \end{tabular}
  \label{tab:liftdrag}
  \end{center}
\end{table}


%%%%%%%%%%%%%%%%%%%%%%%%%%%%%%%%%%%%%%%%%%%%%%%%%%%%%%%%%%%%%%%%%%%%%%%%%%%%%%%%%%%%%%%%%%%%%%%

\section{CITATIONS AND REFERENCES} \label{subsec:ref}

All items included in the References section must be cited in the manuscript, and vice versa. For single-author documents, citations include the Author's last name and publication date, e.g. \cite{Viola2013AMR}. For two Authors' documents, both last names are used, e.g. \cite{Viola2010IJSCTForces}. For more than two Authors, only the first Author's last name is used, followed by et al., e.g. \cite{Viola2021OE}. When the citation is part of the sentence, only the date is between parentheses, e.g. \cite{Viola2013AMR}. Conversely, when the citation is not part of the sentence, then also the Author s name is within parentheses \citep{Viola2013AMR}. Multiple citations that are not part of the sentence should be separated by semicolons \citep{Viola2013AMR,Viola2021OE,Viola2010IJSCTForces}. In case of ambiguity, add a letter to the date such as, for example, in \cite{Viola2013IJHFF} and \cite{Viola2013IJNMF}. 

See the section References for examples of how to include references to journal articles \citep[e.g.][]{Viola2013AMR,Viola2021OE,Viola2010IJSCTForces,Viola2013IJHFF,Viola2013IJNMF}, papers in conference proceedings \citep[e.g.][]{Viola2011CSYS}, book chapters \citep[e.g.]{Viola2013SNAMETrans}, books \citep[e.g.][]{Lanchester1907}, and technical reports \citep[e.g.][]{Carter1984}. When available, include hyperlinks. 

%%%%%%%%%%%%%%%%%%%%%%%%%%%%%%%%%%%%%%%%%%%%%%%%%%%%%%%%%%%%%%%%%%%%%%%%%%%%%%%%%%%%%%%%%%%%%%%

\section{CONCLUSIONS} \label{sec:conclusions}

Each manuscript must include a section named Conclusions, or Discussion and Conclusions, summarising the most important results. 

%%%%%%%%%%%%%%%%%%%%%%%%%%%%%%%%%%%%%%%%%%%%%%%%%%%%%%%%%%%%%%%%%%%%%%%%%%%%%%%%%%%%%%%%%%%%%%%

\section*{ACKNOWLEDGEMENTS}

If the study received funds, this must be noted in the Acknowledgements section after the Conclusions, before the References section or any appendices. The Section Acknowledgements is not numbered. Include both the name of the funding body and, when available, the grant number. Brief personal acknowledgements may also be added. For example, `this work received funds from the European Research Council (grant no. 101001499). The authors are especially grateful to Penelope Smith for her help in undertaking the experiments in the towing tank, and to Olivia Smith for her insightful comments and for the generosity with which she has shared her in-depth knowledge on this subject.'
\end{document}